\documentclass{article}
\usepackage{amsmath,amssymb,xurl}
\usepackage[hidelinks]{hyperref}
% Shared commands for notes in docs/paper-gaps/.

\DeclareUnicodeCharacter{2080}{\ifmmode{}_{0}\else\textsubscript{0}\fi}
\DeclareUnicodeCharacter{2081}{\ifmmode{}_{1}\else\textsubscript{1}\fi}
\DeclareUnicodeCharacter{2082}{\ifmmode{}_{2}\else\textsubscript{2}\fi}
\DeclareUnicodeCharacter{2099}{\ifmmode{}_{n}\else\textsubscript{n}\fi}

\newcommand{\Lean}{\textsc{Lean}}
\ExplSyntaxOn
\NewDocumentCommand{\leanid}{m}{
  \ifmmode
    \textsf{\detokenize{#1}}
  \else
    \group_begin:
    \urlstyle{sf}
    \tl_set:Nn \l_tmpa_tl {#1}
    \tl_replace_all:Nnn \l_tmpa_tl {\_} {_}
    \exp_args:NV \path \l_tmpa_tl
    \group_end:
  \fi
}
\ExplSyntaxOff
\providecommand{\tr}{\operatorname{tr}}
\newcommand{\tnleanIssueBase}{https://github.com/LionSR/TNLean/issues/}
\newcommand{\tnleanPrBase}{https://github.com/LionSR/TNLean/pull/}
\newcommand{\tnleanDocBase}{https://sirui-lu.com/TNLean/docs/}
\newcommand{\ghissue}[1]{\href{\tnleanIssueBase#1}{GitHub issue \##1}}
\newcommand{\ghpr}[1]{\href{\tnleanPrBase#1}{PR \##1}}
\newcommand{\doclink}[1]{\href{\tnleanDocBase#1.html}{\path{#1}}}

% Dirac notation and the identity operator, as in blueprint/src/macros/common.tex.
% \providecommand keeps note-local definitions authoritative.
\usepackage{dsfont}
\providecommand{\ket}[1]{|#1\rangle}
\providecommand{\bra}[1]{\langle #1|}
\providecommand{\braket}[2]{\langle #1 | #2 \rangle}
\providecommand{\ketbra}[2]{|#1\rangle\!\langle #2|}
\providecommand{\Id}{\mathds{1}}
\providecommand{\one}{\mathds{1}}
\providecommand{\C}{\mathbb{C}}
\providecommand{\MN}[1]{M_{#1}(\C)}
\providecommand{\MD}{M_D(\C)}

% Verdict marker: \gapnote{<kind>}{<status>}. Kind is the policy
% classification (clarification, local-correction, scope-restriction,
% unfaithful, false-source, open-gap); status is open, wip (elimination
% underway), resolved, or historical. Machine-read by the site index;
% prints a verdict line.
\newcommand{\gapnote}[2]{\par\noindent\textbf{Verdict:} #1 (#2).\par}

\title{Vacuum Reduced Densities of Closed Flux Strings}
\author{}
\date{}
\begin{document}
\maketitle
\gapnote{scope-restriction}{open}

Schuch, Cirac, and P\'erez-Garc\'ia,
\href{https://arxiv.org/abs/1001.3807v3}{arXiv:1001.3807v3},
Lemma 6.14, source \path{Papers/1001.3807/paper_v3.tex}, lines 2199--2214,
asserts local unobservability of a flux string away from its endpoints.
Theorem 6.7, lines 1995--2025, identifies local physical densities of
commuting torus closure states on a simply connected block.

On a native torus with both periods at least three, let the physical tensor
be regular $G$-isometric for a finite group $G$. Let $R$ have connected
induced occupied graph and simply connected closed-cell realization.
For every finite closed dual path $p:a\to a$ and every $g\in G$, let $M_p$
be its actual inserted physical cut matrix and let $M_0$ use identity
insertions. Then
\begin{align}
    \frac{M_pM_p^\dagger}{\operatorname{tr}(M_pM_p^\dagger)}
    =\frac{M_0M_0^\dagger}{\operatorname{tr}(M_0M_0^\dagger)}.
    \label{eq:closed_flux_vacuum_density}
\end{align}
Both cuts are proved nonzero. Identity plaquette holonomy is derived from
the coincidence of the dual endpoints, then flat-connection classification
and virtual invariance identify the actual state with a commuting seam
closure. The common physical density of those closures gives
(\ref{eq:closed_flux_vacuum_density}). The path may have arbitrary winding;
its global state is not asserted equal to the vacuum.

The declaration is
\leanid{TNLean.PEPS.IsGIsometric.torusClosedFlux_normalized_density_eq_vacuum}.
No factorization, equality of reduced matrices, local parent common-kernel
identification, or nonzero-cut premise is supplied to this theorem.

This proves the closed-string subcase of the density question under the
stated observation geometry. General open flux pairs observed away from
their endpoints still require a separate argument. Endpoint avoidance
alone does not exclude an annular region that detects an enclosed flux.
Neither arbitrary-region path avoidance nor the general open-pair vacuum
reduced-density assertion is claimed; these remain under \ghissue{8703}.
The full-lattice Hamiltonian obligations of \ghissue{8947} are unchanged.
\end{document}
